\chapter*{Liste des symboles}
\begin{twocolumnlist}{.2\textwidth}{.7\textwidth}
    & \large\textbf{Généraux}\\
    $e$ & Charge élémentaire : $1.60217663\times 10^{-19}$ C\\
    $\varepsilon_0$ & Permittivité du vide : $8.8541878188\times 10^{-12}$ F m$^{-1}$\\
    $\hbar$ & Constante de Planck réduite : $1.054571817\times 10^{-34}$ J s\\
    $k_B$ & Constante de Boltzmann : $1.380649\times 10^{-23}$ m$^2$ kg s$^{-2}$ K$^{-1}$\\
    $T$ & Température en kelvins\\
    $\beta$ & Inverse de l'énergie thermique associée à la température : $\beta = 1 /k_BT$\\
    & \\

    & \large\textbf{Réseau cristallin}\\
    $\av_i$ & Vecteurs primitifs du réseau cristallin\\
    $\bv_i$ & Vecteurs primitifs du réseau réciproque du cristal\\
    $\Gv$ & Vecteurs du réseau réciproque du cristal\\
    $\Omega_\text{BvK}$ & Volume du cristal (ou du super-cristal de Born-von-Karman)\\
    $\Omega_\text{uc}$ & Volume de la cellule primitive\\
    $\Omega_\text{sc}$ & Volume de la super-cellule\\
    $N_\text{cells}$ & Nombre de cellules unitaires dans la super-cellule\\
    $N_\text{uc}$ & Nombre de cellules primitives dans le cristal\\ 
    $\Rv_l$ & Vecteurs du réseau cristallin qui identifient la cellule primitive $l$ \\
    $\bm{\tau}_\alpha$ & Position de l'atome $\alpha$ dans une cellule unitaire\\
    & \\

    & \large\textbf{Électrons}\\
    $\kv$ & Vecteur d'onde ou quasi-moment d'un électron\\
    $n$ & Indice de bande d'un électron\\
    $\varepsilon_{n\kv}$ & Énergie d'un électron dans le formalisme de Kohn-Sham\\
    $\ket{\psi_{n\kv}}$ & État quantique de Bloch d'un électron dans le formalisme de Kohn-Sham\\
    $\ket{u_{n\kv}}$ & Partie périodique de l'état de Bloch d'un électron\\
    $\ket{n\kv}$ & État quantique d'une onde plane électronique\\
    $C_{n\kv}(\Gv)$ & Coefficient de l'expansion en ondes planes des états de Bloch électroniques\\
    $\hat{c}^\dag_{n\kv}\,(\hat{c}_{n\kv})$ & Opérateur de création (destruction) d'un électron\\

    & \\

    & \large\textbf{Phonons et photons}\\
    $\qv$ & Vecteur d'onde ou quasi-moment d'un phonon ou d'un photon\\
    $\lambda$ & Polarisation d'un photon\\
    $\nu$ & Branche d'un phonon\\
    $\mathbf{e}_{\lambda\qv}$ & Vecteur de polarisation d'un photon\\
    $\mathbf{e}_{\nu\qv}$ & Vecteur de polarisation d'un phonon \\
    $\omega_{\qv}$ & Fréquence d'un photon\\
    $\omega_{\nu\qv}$ & Fréquence d'un phonon\\
    $\hat{a}_{\lambda\qv}^\dag\,(\hat{a}_{\lambda\qv})$ & Opérateur de création (destruction) d'un photon\\
    $\hat{b}_{\nu\qv}^\dag\,(\hat{b}_{\nu\qv})$ & Opérateur de création (destruction) d'un phonon\\
    $\hat{A}_{\nu\qv}$ & Opérateur proportionnel aux déplacements d'un phonon : $\hat{A}_{\nu\qv}=\hat{b}_{\nu\qv}+\hat{b}^\dag_{\nu-\qv}$\\
    & \\

    & \large\textbf{Couplage}\\
    $g_{nm\lambda}(\kv,\qv)$ & Élément de matrice de couplage électron-photon\\
    $g_{nm\nu}(\kv,\qv)$ & Élément de matrice de couplage électron-phonon\\
    $M_{nm}(\kv,\kv')$ & Élément de matrice de couplage électron-défaut\\
    & \\

    & \large\textbf{Fonctions de Green et \textit{self}-énergie} \\
    $\mathcal{D}_{\nu\qv}(i\nu_j)$ & Fonction de Green de Matsubara phononique aux fréquences de Matsubara\\
    $\mathcal{D}^{(0)}_{\nu\qv}(i\nu_j)$ & Fonction de Green de Matsubara phononique libre aux fréquences de Matsubara\\
    $\mathcal{G}_{n\kv}(i\omega_j)$ & Fonction de Green de Matsubara électronique aux fréquences de Matsubara\\
    $\mathcal{G}^\text{ed}_{n\kv}(i\omega_j)$ & Fonction de Green de Matsubara électronique due aux défauts aux fréquences de Matsubara\\
    $\mathcal{G}^\text{eph}_{n\kv}(i\omega_j)$ & Fonction de Green de Matsubara électronique due phonons aux fréquences de Matsubara\\
    $\mathcal{G}^{(0)}_{n\kv}(i\omega_j)$ & Fonction de Green de Matsubara électronique libre aux fréquences de Matsubara\\
    $\bar{\mathcal{G}}^\text{ed}_{n\kv}(i\omega_j)$ & Fonction de Green de Matsubara électronique moyenée pour une collections de défauts aux fréquences de Matsubara\\
    $i\omega_j$ & Fréquence de Matsubara fermionique\\
    $i\nu_j$ & Fréquence de Matsubara bosonique\\
    $\Pi_{\nu\qv}(\omega)$ & \textit{Self}-énergie phononique due aux interactions avec les électrons\\
    $\Sigma^\text{ed}_{n\kv}(\omega)$ & \textit{Self}-énergie électronique due aux interactions avec les défauts\\
    $\Sigma^\text{eph}_{n\kv}(\omega)$ & \textit{Self}-énergie électronique due aux interactions avec les phonons\\
    $S(\beta, 0)$ & Matrice $S$ dans le formalisme de Matsubara\\
    $T_{mn}(\kv',\kv)$ & Matrice $T$ pour le couplage des électrons avec un seul défaut\\

    & \\
    & \large\textbf{Hamiltoniens}\\
    $\hat{H}$ & Hamiltonien général\\
    $\hat{H}_0$ & Hamiltonien libre général\\
    $\hat{H}_\text{e}$ & Hamiltonien électronique\\
    $\hat{H}_\text{ed}$ & Hamiltonien de couplage électron-défaut\\
    $\hat{H}_\text{eem}$ & Hamiltonien de couplage électron-photon\\
    $\hat{H}_\text{em}$ & Hamiltonien des photons libres\\  
    $\hat{H}_\text{eph}$ & Hamiltonien de couplage électron-phonon\\
    $\hat{H}_\text{i}$ & Hamiltonien ionique\\
    $\hat{H}_\text{int}$ & Hamiltonien d'interaction général\\
    $\hat{H}_\text{KS}$ & Hamiltonien de Kohn-Sham dans le formalisme de la théorie de la fonctionnelle de la densité\\
    $\hat{H}_\text{ph}$ & Hamiltonien des phonons libres\\
\end{twocolumnlist}